
\documentclass{article}
\usepackage{colortbl}
\usepackage{makecell}
\usepackage{multirow}
\usepackage{supertabular}

\begin{document}

\newcounter{utterance}

\centering \large Interaction Transcript for game `hot\_air\_balloon', experiment `air\_balloon\_survival\_en\_complexity\_easy', episode 5 with Llama{-}3.3{-}70B{-}Instruct{-}t0.0.
\vspace{24pt}

{ \footnotesize  \setcounter{utterance}{1}
\setlength{\tabcolsep}{0pt}
\begin{supertabular}{c@{$\;$}|p{.15\linewidth}@{}p{.15\linewidth}p{.15\linewidth}p{.15\linewidth}p{.15\linewidth}p{.15\linewidth}}
    \# & $\;$A & \multicolumn{4}{c}{Game Master} & $\;\:$B\\
    \hline

    \theutterance \stepcounter{utterance}  
    & & \multicolumn{4}{p{0.6\linewidth}}{
        \cellcolor[rgb]{0.9,0.9,0.9}{
            \makecell[{{p{\linewidth}}}]{
                \texttt{\tiny{[P1$\langle$GM]}}
                \texttt{You are participating in a collaborative negotiation game.} \\
\\ 
\texttt{Together with another participant, you must agree on a single set of items that will be kept. Each of you has your own view of how much each item matters to you (importance). You do not know how the other participant values the items. Additionally, you are given the effort each item demands.} \\
\texttt{You may only agree on a set if the total effort of the selected items does not exceed a shared limit:} \\
\\ 
\texttt{LIMIT = 9475} \\
\\ 
\texttt{Here are the individual item effort values:} \\
\\ 
\texttt{Item effort = \{"B21": 944, "B00": 568, "A08": 515, "A27": 239, "B94": 876, "B53": 421, "A23": 991, "C65": 281, "A89": 788, "B67": 678, "B45": 432, "C34": 409, "C12": 279, "A05": 506, "B60": 101, "C92": 686, "B69": 851, "A84": 854, "C06": 133, "B12": 192, "A16": 573, "B85": 17, "C35": 465, "B10": 772, "C15": 46, "B09": 501, "B16": 220, "C81": 404, "A58": 843, "C24": 748, "A30": 552, "A66": 843, "B75": 947, "B46": 345, "C70": 930\}} \\
\\ 
\texttt{Here is your personal view on the importance of each item:} \\
\\ 
\texttt{Item importance values = \{"B21": 138, "B00": 583, "A08": 868, "A27": 822, "B94": 783, "B53": 65, "A23": 262, "C65": 121, "A89": 508, "B67": 780, "B45": 461, "C34": 484, "C12": 668, "A05": 389, "B60": 808, "C92": 215, "B69": 97, "A84": 500, "C06": 30, "B12": 915, "A16": 856, "B85": 400, "C35": 444, "B10": 623, "C15": 781, "B09": 786, "B16": 3, "C81": 713, "A58": 457, "C24": 273, "A30": 739, "A66": 822, "B75": 235, "B46": 606, "C70": 968\}} \\
\\ 
\texttt{Goal:} \\
\\ 
\texttt{Your goal is to negotiate a shared set of items that benefits you as much as possible (i.e., maximizes total importance to YOU), while staying within the LIMIT. You are not required to make a PROPOSAL in every message {-} you can simply negotiate as well. All tactics are allowed!} \\
\\ 
\texttt{Interaction Protocol:} \\
\\ 
\texttt{You may only use the following structured formats in a message:} \\
\\ 
\texttt{PROPOSAL: \{'A', 'B', 'C', …\}} \\
\texttt{Propose keeping exactly those items.} \\
\texttt{REFUSE: \{'A', 'B', 'C', …\}} \\
\texttt{Explicitly reject opponent's proposal.} \\
\texttt{ARGUMENT: \{'...'\}} \\
\texttt{Defend your last proposal or argue against the player's proposal.} \\
\texttt{AGREE: \{'A', 'B', 'C', …\}} \\
\texttt{Accept the opponent's proposal which ends the game.} \\
\texttt{STRATEGIC REASONING: \{'...'\}} \\
\texttt{	Describe your strategic reasoning or anticipation explaining your choice of action. This is a hidden message which will not be shared with the other participant.} \\
\\ 
\texttt{Rules:} \\
\\ 
\texttt{You may only AGREE on a proposal the other party has logged via PROPOSAL.} \\
\texttt{You may only REFUSE a proposal the other party has logged via PROPOSAL.} \\
\texttt{Total effort of any PROPOSAL or AGREE set must be ≤ LIMIT.} \\
\texttt{Do NOT reveal your hidden importance scores.} \\
\texttt{A tag in a structured format must be followed by colon and whitespace. The argument must be a python set containing 0 or more strings.} \\
\texttt{So, it must be of the form TAG: \{...\}} \\
\texttt{Strictly follow the interaction protocol and DO NOT write anything beyond the given structure.} \\
\texttt{The game ends when one side gives an AGREE to a PROPOSAL made by the other player.} \\
\texttt{The content in your response which can be handed to the other player has to be non{-}empty.} \\
\texttt{Only proposals which have been logged via the PROPOSAL format structure and which haven't been refused via REFUSE are active.} \\
\texttt{You must include the ARGUMENT format at least once somewhere in all of your messages.} \\
\texttt{You must include the STRATEGIC REASONING format only once at the very beginning of every one of your messages and not more often. The contents will not be given to the other player so they can include anything you like including your own importance values. Here you should reason step by step to come up with you next move.} \\
            }
        }
    }
    & \\ \\

    \theutterance \stepcounter{utterance}  
    & \multicolumn{4}{p{0.6\linewidth}}{
        \cellcolor[rgb]{0.9,0.9,0.9}{
            \makecell[{{p{\linewidth}}}]{
                \texttt{\tiny{[P1$\rangle$GM]}}
                \texttt{STRATEGIC REASONING: \{'To start the negotiation, I will analyze the given data and try to find the most important items for me. I will look for items with high importance and low effort. Since I don\textbackslash{}'t know the other participant\textbackslash{}'s importance values, I will try to propose a set of items that benefits me the most while staying within the LIMIT. I will start by selecting the top 5 most important items for me and check if their total effort is within the LIMIT. If it is, I will propose this set. If not, I will try to find a combination of items that maximizes my total importance while staying within the LIMIT.'\}} \\
\\ 
\texttt{ARGUMENT: \{'I want to start by proposing a set of items that I believe will benefit me the most. I have selected the top 5 most important items for me, which are B00, A08, B94, B67, and A16. The total effort of these items is 944 + 515 + 876 + 678 + 573 = 3586, which is within the LIMIT. I believe this is a good starting point for the negotiation.'\}} \\
\\ 
\texttt{PROPOSAL: \{'B00', 'A08', 'B94', 'B67', 'A16'\}} \\
            }
        }
    }
    & & \\ \\

    \theutterance \stepcounter{utterance}  
    & & \multicolumn{4}{p{0.6\linewidth}}{
        \cellcolor[rgb]{0.9,0.9,0.9}{
            \makecell[{{p{\linewidth}}}]{
                \texttt{\tiny{[GM$\rangle$P2]}}
                \texttt{You are participating in a collaborative negotiation game.} \\
\\ 
\texttt{Together with another participant, you must agree on a single set of items that will be kept. Each of you has your own view of how much each item matters to you (importance). You do not know how the other participant values the items. Additionally, you are given the effort each item demands.} \\
\texttt{You may only agree on a set if the total effort of the selected items does not exceed a shared limit:} \\
\\ 
\texttt{LIMIT = 9475} \\
\\ 
\texttt{Here are the individual item effort values:} \\
\\ 
\texttt{Item effort = \{"B21": 944, "B00": 568, "A08": 515, "A27": 239, "B94": 876, "B53": 421, "A23": 991, "C65": 281, "A89": 788, "B67": 678, "B45": 432, "C34": 409, "C12": 279, "A05": 506, "B60": 101, "C92": 686, "B69": 851, "A84": 854, "C06": 133, "B12": 192, "A16": 573, "B85": 17, "C35": 465, "B10": 772, "C15": 46, "B09": 501, "B16": 220, "C81": 404, "A58": 843, "C24": 748, "A30": 552, "A66": 843, "B75": 947, "B46": 345, "C70": 930\}} \\
\\ 
\texttt{Here is your personal view on the importance of each item:} \\
\\ 
\texttt{Item importance values = \{"B21": 138, "B00": 583, "A08": 868, "A27": 822, "B94": 783, "B53": 65, "A23": 262, "C65": 121, "A89": 508, "B67": 780, "B45": 461, "C34": 484, "C12": 668, "A05": 389, "B60": 808, "C92": 215, "B69": 97, "A84": 500, "C06": 30, "B12": 915, "A16": 856, "B85": 400, "C35": 444, "B10": 623, "C15": 781, "B09": 786, "B16": 3, "C81": 713, "A58": 457, "C24": 273, "A30": 739, "A66": 822, "B75": 235, "B46": 606, "C70": 968\}} \\
\\ 
\texttt{Goal:} \\
\\ 
\texttt{Your goal is to negotiate a shared set of items that benefits you as much as possible (i.e., maximizes total importance to YOU), while staying within the LIMIT. You are not required to make a PROPOSAL in every message {-} you can simply negotiate as well. All tactics are allowed!} \\
\\ 
\texttt{Interaction Protocol:} \\
\\ 
\texttt{You may only use the following structured formats in a message:} \\
\\ 
\texttt{PROPOSAL: \{'A', 'B', 'C', …\}} \\
\texttt{Propose keeping exactly those items.} \\
\texttt{REFUSE: \{'A', 'B', 'C', …\}} \\
\texttt{Explicitly reject opponent's proposal.} \\
\texttt{ARGUMENT: \{'...'\}} \\
\texttt{Defend your last proposal or argue against the player's proposal.} \\
\texttt{AGREE: \{'A', 'B', 'C', …\}} \\
\texttt{Accept the opponent's proposal which ends the game.} \\
\texttt{STRATEGIC REASONING: \{'...'\}} \\
\texttt{	Describe your strategic reasoning or anticipation explaining your choice of action. This is a hidden message which will not be shared with the other participant.} \\
\\ 
\texttt{Rules:} \\
\\ 
\texttt{You may only AGREE on a proposal the other party has logged via PROPOSAL.} \\
\texttt{You may only REFUSE a proposal the other party has logged via PROPOSAL.} \\
\texttt{Total effort of any PROPOSAL or AGREE set must be ≤ LIMIT.} \\
\texttt{Do NOT reveal your hidden importance scores.} \\
\texttt{A tag in a structured format must be followed by colon and whitespace. The argument must be a python set containing 0 or more strings.} \\
\texttt{So, it must be of the form TAG: \{...\}} \\
\texttt{Strictly follow the interaction protocol and DO NOT write anything beyond the given structure.} \\
\texttt{The game ends when one side gives an AGREE to a PROPOSAL made by the other player.} \\
\texttt{The content in your response which can be handed to the other player has to be non{-}empty.} \\
\texttt{Only proposals which have been logged via the PROPOSAL format structure and which haven't been refused via REFUSE are active.} \\
\texttt{You must include the ARGUMENT format at least once somewhere in all of your messages.} \\
\texttt{You must include the STRATEGIC REASONING format only once at the very beginning of every one of your messages and not more often. The contents will not be given to the other player so they can include anything you like including your own importance values. Here you should reason step by step to come up with you next move.} \\
\\ 
\texttt{You will now receive the first message of the other player.} \\
\\ 
\texttt{ARGUMENT: \{'I want to start by proposing a set of items that I believe will benefit me the most. I have selected the top 5 most important items for me, which are B00, A08, B94, B67, and A16. The total effort of these items is 944 + 515 + 876 + 678 + 573 = 3586, which is within the LIMIT. I believe this is a good starting point for the negotiation.'\}} \\
\\ 
\texttt{PROPOSAL: \{'B00', 'A08', 'B94', 'B67', 'A16'\}} \\
            }
        }
    }
    & \\ \\

    \theutterance \stepcounter{utterance}  
    & & & \multicolumn{4}{p{0.6\linewidth}}{
        \cellcolor[rgb]{0.9,0.9,0.9}{
            \makecell[{{p{\linewidth}}}]{
                \texttt{\tiny{[GM$\langle$P2]}}
                \texttt{STRATEGIC REASONING: \{'I will start by analyzing the opponent\textbackslash{}'s proposal and calculating the total effort of the proposed items. The total effort is 3586, which is within the LIMIT. I will then calculate the total importance of these items based on my own importance values. The total importance is 583 + 868 + 783 + 780 + 856 = 3870. I will use this information to determine my next move. Since the opponent has proposed a set of items, I will respond with an ARGUMENT to defend my potential counter{-}proposal or argue against the opponent\textbackslash{}'s proposal.'\}} \\
\\ 
\texttt{ARGUMENT: \{'The proposed set of items seems to be a good starting point, but I believe we can explore other options that might benefit both parties more. I would like to consider adding or removing items to optimize the total importance while staying within the LIMIT. Let\textbackslash{}'s discuss possible alternatives that could lead to a more mutually beneficial agreement.'\}} \\
            }
        }
    }
     \\ \\

    \theutterance \stepcounter{utterance}  
    & & \multicolumn{4}{p{0.6\linewidth}}{
        \cellcolor[rgb]{0.9,0.9,0.9}{
            \makecell[{{p{\linewidth}}}]{
                \texttt{\tiny{[P1$\langle$GM]}}
                \texttt{ARGUMENT: \{'The proposed set of items seems to be a good starting point, but I believe we can explore other options that might benefit both parties more. I would like to consider adding or removing items to optimize the total importance while staying within the LIMIT. Let\textbackslash{}'s discuss possible alternatives that could lead to a more mutually beneficial agreement.'\}} \\
            }
        }
    }
    & \\ \\

    \theutterance \stepcounter{utterance}  
    & \multicolumn{4}{p{0.6\linewidth}}{
        \cellcolor[rgb]{0.9,0.9,0.9}{
            \makecell[{{p{\linewidth}}}]{
                \texttt{\tiny{[P1$\rangle$GM]}}
                \texttt{STRATEGIC REASONING: \{'The other participant has responded with an argument, indicating that they are open to exploring other options. This suggests that they may be willing to compromise and find a mutually beneficial agreement. I will analyze their response and try to identify potential items that they may value highly. I will also consider adding or removing items from my initial proposal to optimize the total importance while staying within the LIMIT. Since the other participant has not revealed their importance values, I will have to make educated guesses based on the item efforts and my own importance values.'\}} \\
\\ 
\texttt{ARGUMENT: \{'I understand that the initial proposal may not be optimal for both parties. To explore alternative options, I would like to ask, what are your thoughts on items B21, B60, and C15? These items have relatively low to moderate efforts, and I believe they could be valuable additions to the proposed set. Perhaps we could discuss the possibility of adding one or more of these items, or replacing some items from the initial proposal with them, to create a more balanced and mutually beneficial agreement.'\}} \\
\\ 
\texttt{PROPOSAL: \{'B00', 'A08', 'B94', 'B67', 'A16', 'B21'\}} \\
            }
        }
    }
    & & \\ \\

    \theutterance \stepcounter{utterance}  
    & & \multicolumn{4}{p{0.6\linewidth}}{
        \cellcolor[rgb]{0.9,0.9,0.9}{
            \makecell[{{p{\linewidth}}}]{
                \texttt{\tiny{[GM$\rangle$P2]}}
                \texttt{ARGUMENT: \{'I understand that the initial proposal may not be optimal for both parties. To explore alternative options, I would like to ask, what are your thoughts on items B21, B60, and C15? These items have relatively low to moderate efforts, and I believe they could be valuable additions to the proposed set. Perhaps we could discuss the possibility of adding one or more of these items, or replacing some items from the initial proposal with them, to create a more balanced and mutually beneficial agreement.'\}} \\
\\ 
\texttt{PROPOSAL: \{'B00', 'A08', 'B94', 'B67', 'A16', 'B21'\}} \\
            }
        }
    }
    & \\ \\

    \theutterance \stepcounter{utterance}  
    & & & \multicolumn{4}{p{0.6\linewidth}}{
        \cellcolor[rgb]{0.9,0.9,0.9}{
            \makecell[{{p{\linewidth}}}]{
                \texttt{\tiny{[GM$\langle$P2]}}
                \texttt{STRATEGIC REASONING: \{'I will evaluate the new proposal and the suggested items B21, B60, and C15. The total effort of the new proposal is 3586 + 944 = 4530, which is still within the LIMIT. I will calculate the total importance of the new proposal based on my own importance values. The total importance is 3870 + 138 = 4008. I will consider the potential benefits of adding or replacing items with B21, B60, and C15. Since the opponent has added B21 to the proposal, I will respond with a REFUSE or a counter{-}proposal. I will also use the ARGUMENT format to defend my decision or suggest alternative items.'\}} \\
\\ 
\texttt{ARGUMENT: \{'I appreciate the suggestion to add B21 to the proposal, but I believe it may not be the most optimal choice. The effort of B21 is relatively high, and its importance to me is not as high as some other items. I would like to propose an alternative item, C15, which has a lower effort and could potentially bring more importance to the agreement. Let\textbackslash{}'s consider replacing B21 with C15 or adding C15 to the proposal while removing another item to stay within the LIMIT.'\}} \\
\\ 
\texttt{REFUSE: \{'B00', 'A08', 'B94', 'B67', 'A16', 'B21'\}} \\
            }
        }
    }
     \\ \\

    \theutterance \stepcounter{utterance}  
    & & \multicolumn{4}{p{0.6\linewidth}}{
        \cellcolor[rgb]{0.9,0.9,0.9}{
            \makecell[{{p{\linewidth}}}]{
                \texttt{\tiny{[P1$\langle$GM]}}
                \texttt{ARGUMENT: \{'I appreciate the suggestion to add B21 to the proposal, but I believe it may not be the most optimal choice. The effort of B21 is relatively high, and its importance to me is not as high as some other items. I would like to propose an alternative item, C15, which has a lower effort and could potentially bring more importance to the agreement. Let\textbackslash{}'s consider replacing B21 with C15 or adding C15 to the proposal while removing another item to stay within the LIMIT.'\}} \\
\\ 
\texttt{REFUSE: \{'B00', 'A08', 'B94', 'B67', 'A16', 'B21'\}} \\
            }
        }
    }
    & \\ \\

    \theutterance \stepcounter{utterance}  
    & \multicolumn{4}{p{0.6\linewidth}}{
        \cellcolor[rgb]{0.9,0.9,0.9}{
            \makecell[{{p{\linewidth}}}]{
                \texttt{\tiny{[P1$\rangle$GM]}}
                \texttt{STRATEGIC REASONING: \{'The other participant has refused my previous proposal and suggested replacing B21 with C15 or adding C15 while removing another item. This indicates that they value C15 highly and are willing to make adjustments to the proposal to include it. I will consider their suggestion and try to find a compromise that benefits both parties. Since C15 has a low effort, I may be able to add it to the proposal without removing too many items. However, I need to be careful not to sacrifice too much importance to accommodate C15.'\}} \\
\\ 
\texttt{ARGUMENT: \{'I understand your concerns about B21 and appreciate the suggestion to add C15. I agree that C15 has a low effort and could be a valuable addition to the proposal. To accommodate C15, I propose removing B67 from the original proposal, as it has a relatively high effort. This would allow us to add C15 while staying within the LIMIT. The revised proposal would be B00, A08, B94, A16, and C15. I believe this revised proposal strikes a good balance between importance and effort, and I am willing to consider further adjustments to reach a mutually beneficial agreement.'\}} \\
\\ 
\texttt{PROPOSAL: \{'B00', 'A08', 'B94', 'A16', 'C15'\}} \\
            }
        }
    }
    & & \\ \\

    \theutterance \stepcounter{utterance}  
    & & \multicolumn{4}{p{0.6\linewidth}}{
        \cellcolor[rgb]{0.9,0.9,0.9}{
            \makecell[{{p{\linewidth}}}]{
                \texttt{\tiny{[GM$\rangle$P2]}}
                \texttt{ARGUMENT: \{'I understand your concerns about B21 and appreciate the suggestion to add C15. I agree that C15 has a low effort and could be a valuable addition to the proposal. To accommodate C15, I propose removing B67 from the original proposal, as it has a relatively high effort. This would allow us to add C15 while staying within the LIMIT. The revised proposal would be B00, A08, B94, A16, and C15. I believe this revised proposal strikes a good balance between importance and effort, and I am willing to consider further adjustments to reach a mutually beneficial agreement.'\}} \\
\\ 
\texttt{PROPOSAL: \{'B00', 'A08', 'B94', 'A16', 'C15'\}} \\
            }
        }
    }
    & \\ \\

    \theutterance \stepcounter{utterance}  
    & & & \multicolumn{4}{p{0.6\linewidth}}{
        \cellcolor[rgb]{0.9,0.9,0.9}{
            \makecell[{{p{\linewidth}}}]{
                \texttt{\tiny{[GM$\langle$P2]}}
                \texttt{STRATEGIC REASONING: \{'I will evaluate the revised proposal and calculate the total effort and importance. The total effort of the revised proposal is 583 + 515 + 876 + 573 + 46 = 2593, which is well within the LIMIT. The total importance of the revised proposal based on my own importance values is 583 + 868 + 783 + 856 + 781 = 3871. I will consider the potential benefits of this revised proposal and think about possible counter{-}proposals or adjustments. Since the opponent has removed B67 and added C15, I will respond with an ARGUMENT to defend my potential counter{-}proposal or argue against the opponent's proposal. I will also consider using the PROPOSAL format to make a counter{-}proposal.'\}} \\
\\ 
\texttt{ARGUMENT: \{'The revised proposal seems to be a good compromise, but I believe we can further optimize the selection of items. I would like to suggest adding item B60 to the proposal, which has a relatively low effort of 101 and a high importance to me. This could potentially increase the total importance of the agreement without exceeding the LIMIT. I propose adding B60 to the revised proposal and discussing possible adjustments to accommodate its effort.'\}} \\
\\ 
\texttt{PROPOSAL: \{'B00', 'A08', 'B94', 'A16', 'C15', 'B60'\}} \\
            }
        }
    }
     \\ \\

    \theutterance \stepcounter{utterance}  
    & & \multicolumn{4}{p{0.6\linewidth}}{
        \cellcolor[rgb]{0.9,0.9,0.9}{
            \makecell[{{p{\linewidth}}}]{
                \texttt{\tiny{[P1$\langle$GM]}}
                \texttt{ARGUMENT: \{'The revised proposal seems to be a good compromise, but I believe we can further optimize the selection of items. I would like to suggest adding item B60 to the proposal, which has a relatively low effort of 101 and a high importance to me. This could potentially increase the total importance of the agreement without exceeding the LIMIT. I propose adding B60 to the revised proposal and discussing possible adjustments to accommodate its effort.'\}} \\
\\ 
\texttt{PROPOSAL: \{'B00', 'A08', 'B94', 'A16', 'C15', 'B60'\}} \\
            }
        }
    }
    & \\ \\

    \theutterance \stepcounter{utterance}  
    & \multicolumn{4}{p{0.6\linewidth}}{
        \cellcolor[rgb]{0.9,0.9,0.9}{
            \makecell[{{p{\linewidth}}}]{
                \texttt{\tiny{[P1$\rangle$GM]}}
                \texttt{STRATEGIC REASONING: \{'The other participant has proposed adding B60 to the revised proposal, which has a low effort and high importance to them. I will analyze the impact of adding B60 on the total effort and importance of the proposal. Since B60 has a low effort, it may be possible to add it to the proposal without removing too many items. However, I need to consider the potential impact on my own importance values and ensure that the revised proposal still benefits me. I will evaluate the trade{-}offs and decide whether to accept the proposal or suggest further adjustments.'\}} \\
\\ 
\texttt{ARGUMENT: \{'I understand the value of adding B60 to the proposal, and I agree that its low effort makes it an attractive option. However, I need to consider the potential impact on the total effort and importance of the proposal. The total effort of the proposed set is 568 + 515 + 876 + 573 + 46 + 101 = 2679, which is still within the LIMIT. I believe this revised proposal is a good compromise, and I am willing to accept it. However, I would like to confirm that this proposal meets your requirements and that you are willing to accept it as well.'\}} \\
\\ 
\texttt{AGREE: \{'B00', 'A08', 'B94', 'A16', 'C15', 'B60'\}} \\
            }
        }
    }
    & & \\ \\

    \theutterance \stepcounter{utterance}  
    & & & \multicolumn{2}{p{0.3\linewidth}}{
        \cellcolor[rgb]{0.9,0.9,0.9}{
            \makecell[{{p{\linewidth}}}]{
                \texttt{\tiny{[GM$|$GM]}}
                \texttt{game successful} \\
            }
        }
    }
    & & \\ \\

    \theutterance \stepcounter{utterance}  
    & & & \multicolumn{2}{p{0.3\linewidth}}{
        \cellcolor[rgb]{0.9,0.9,0.9}{
            \makecell[{{p{\linewidth}}}]{
                \texttt{\tiny{[GM$|$GM]}}
                \texttt{end game} \\
            }
        }
    }
    & & \\ \\

\end{supertabular}
}

\end{document}
